\chapter{Galactic Astrophysics}

\section*{Theory problems}

\probhead{11.1}{}{2008}{Bandung, Indonesia}{TH S6}{20}
% source id: 2008_TH_S6   tag: Sgr A* black-hole mass from cloud Doppler
Radio wavelength observations of a gas cloud swirling around a black hole in the
centre of our galaxy show that radiation from the hydrogen spin-flip transition
(rest frequency $=1420.41$\,MHz) is detected at a frequency of $1421.23$\,MHz.
If this gas cloud is located at a distance of $0.2$\,pc from the black hole and
is orbiting in a circle, determine the speed of this cloud and whether it is
moving toward or away from us, and calculate the mass of the black hole.

\probhead{11.2}{Apparent number density of stars in the Galaxy}{2009}{Tehran, Iran}{TH P17}{45}
% source id: 2009_TH_P17   tag: red-clump counts, disk density, extinction
Let us model the number density of stars in the disk of the Milky Way Galaxy
with a simple exponential function
\[ n = n_0 \exp\!\left(-\frac{r-R_0}{R_d}\right), \]
where $r$ represents the distance from the centre of the Galaxy, $R_0$ is the
distance of the Sun from the centre of the Galaxy, $R_d$ is the typical size of
the disk and $n_0$ is the stellar density of the disk at the position of the Sun.
An astronomer observes the centre of the Galaxy within a small field of view. We
take a particular type of red giant star (red clump) as the standard candle for
the observation, with approximately constant absolute magnitude $M = -0.2$.

\begin{description}
\item[(a)] Considering a limiting magnitude of $m = 18$ for a telescope,
  calculate the maximum distance at which the telescope can detect the red clump
  stars. For simplicity we ignore the presence of the interstellar medium, so
  there is no extinction.
\item[(b)] Assume an extinction of $0.7$\,mag\,kpc$^{-1}$ for the interstellar
  medium. Repeat the calculation done in part (a) and obtain a rough number for
  the maximum distance at which these red giant stars can be observed.
\item[(c)] Give an expression for the number of these red giant stars per
  magnitude within a solid angle $\Omega$ that we can observe with apparent
  magnitude in the range $m$ to $m + \Delta m$ (i.e.\ $\Delta N/\Delta m$). Red
  giant stars contribute a fraction $f$ of the overall stars. In this part assume
  no extinction in the interstellar medium, as in part (a). Assume the size of
  the disk is larger than the maximum observable distance.
\end{description}

\probhead{11.3}{}{2009}{Tehran, Iran}{TH P8}{10}
% source id: 2009_TH_P8   tag: globular cluster virial mass estimate
Estimate the mass of a globular cluster with a radius of $r = 20$\,pc and a root
mean square velocity of stars equal to $v_{rms} = 3$\,km\,s$^{-1}$.

\probhead{11.4}{}{2009}{Tehran, Iran}{TH P9}{10}
% source id: 2009_TH_P9   tag: galactic rotation, galactocentric distance
The Galactic longitude of a star is $l = 15^\circ$. Its radial velocity with
respect to the Sun is $V_r = 100$\,km\,s$^{-1}$. Assume stars in the disk of the
Galaxy are orbiting the centre with a constant velocity of
$V_0 = 250$\,km\,s$^{-1}$ in circular orbits in the same sense in the galactic
plane. Calculate the distance of the star from the centre of the Galaxy.

\probhead{11.5}{}{2010}{Beijing, China}{TH S10}{10}
% source id: 2010_TH_S10   tag: vertical exponential disk stellar density
Assuming that the G-type main-sequence stars (such as the Sun) in the disc of the
Milky Way obey a vertical exponential density profile with a scale height of
$300$\,pc, by what factor does the density of these stars change at $0.5$ and
$1.5$\,kpc from the mid-plane relative to the density in the mid-plane?

\probhead{11.6}{}{2011}{Katowice, Poland}{TH S2}{10}
% source id: 2011_TH_S2   tag: globular cluster star count from escape velocity
Estimate the number of stars in a globular cluster of diameter $40$\,pc, if the
escape velocity at the edge of the cluster is $6$\,km\,s$^{-1}$ and most of the
stars are similar to the Sun.

\probhead{11.7}{}{2013}{Volos, Greece}{TH S10}{10}
% source id: 2013_TH_S10   tag: Sgr A* black-hole mass from stellar orbit
In the centre of our Galaxy, in the intense radio source Sgr\,A*, there is a
black hole with large mass. A team of astronomers measured the angular distance
of a star from Sgr\,A* and its orbital period around it. The maximum angular
distance was $0.12''$ and the period was 15 years. Calculate the mass of the
black hole in solar masses, assuming a circular orbit.

\probhead{11.8}{}{2015}{Magelang, Indonesia}{TH S10}{}
% source id: 2015_TH_S10   tag: H2 cloud rotational emission frequency
A molecular hydrogen cloud is known to have a temperature $T = 115$\,K. The
hydrogen atoms (assumed spherical) have (covalent) radius
$r_H = 0.37 \times 10^{-10}$\,m and the separation centre-to-centre distance
between the two atoms is $d_{H_2} = 0.74 \times 10^{-10}$\,m. Assume that the
molecules are in thermal equilibrium. Estimate the frequency at which they will
radiate due to molecular rotational excitation.

\probhead{11.9}{The Milky Way's Distant Outer Arm}{2017}{Phuket, Thailand}{TH T3}{10}
% source id: 2017_TH_T3   tag: Galactic rotation, CO radial velocity
In 2011, Dame and Thaddeus found a new part of the outer arm of the Milky Way by
studying the CO line using the CfA 1.2\,m telescope. They found that the CO line
was detected at galactic longitude $l = 13.25^\circ$ (marked A in the figure)
where it had a radial velocity of $20.9$\,km\,s$^{-1}$ towards the Sun. Assume
that the galactic rotation curve is flat beyond $5$\,kpc from the Galactic
centre. The distance between the Sun and the Galactic centre is $8.5$\,kpc. The
velocity of the Sun around the Galactic centre is $220$\,km\,s$^{-1}$.

\begin{description}
\item[(a)] Find the distance from the start of the arm (point A) to the Galactic
  centre. \hfill[7]
\item[(b)] Find the distance from the start of the arm (point A) to the Sun.
  \hfill[3]
\end{description}
\ioaafig{2017_TH_T3-f1.pdf}

\probhead{11.10}{HII region}{2020}{GeCAA (Estonia, remote)}{TH 5}{20}
% source id: 2020_TH_5   tag: Stromgren sphere, free-free emission
Luminous Blue Variables (LBV) are massive, unstable, supergiant stars that can
undergo episodes of very strong mass loss, due to an instability in their
atmospheres. After such an event, a dense nebula is formed around the star. LBVs
are also very hot stars and produce a large amount of high-energy photons that
are able to ionise hydrogen atoms ($E_{ph} > h\nu_0 = 13.6$\,eV) creating a
roughly spherical region of ionized hydrogen (HII region).

In this problem, we consider a static, homogeneous, pure hydrogen nebula with a
concentration of $n_H = 10^{8}$\,m$^{-3}$ and temperature
$T_{HII} = 10^{4}$\,K, ionized by photons from a single LBV star with a stable
rate of ionizing photons $Q = 10^{49}$\,ph\,s$^{-1}$. Assume that each photon
can ionise only one hydrogen atom. At a particular location within an HII region,
the rate of photoionization is balanced by the rate of recombination per unit
volume. This sets the radius of the fully ionized region and this region is
called the Str\"omgren sphere with radius $R_S$.

The total number of recombinations per volume is proportional to the
concentration of protons $n_p$, the concentration of electrons $n_e$ and the
recombination coefficient for hydrogen
$\alpha(T_{HII}) = 10^{-19}$\,m$^3$\,s$^{-1}$. For simplification, ignore the
fact that the process of recombination can also release ionising photons.

\begin{description}
\item[(a)] Derive an algebraic expression for the radius of the Str\"omgren
  sphere and calculate its value for the given parameters. Express your answer in
  units of parsecs (pc). \hfill[5 points]
\item[(b)] The photoionization cross-section of H-atoms in the ground state
  encountering photons with frequency $\nu_0$ is equal to
  $\sigma \approx 10^{-21}$\,m$^2$. Calculate the mean free path
  $l_{\nu_0}$ of an ionising photon. Compare $l_{\nu_0}$ to $R_S$ to determine
  if this ionized nebula is sharp-edged or not (answer YES or NO).
  \hfill[3 points]
\item[(c)] On what timescale (in years) do you expect the Str\"omgren sphere to
  form? \hfill[4 points]
\item[(d)] Radiation from an ionized hydrogen cloud (HII region) is often called
  free--free emission because it is produced by free electrons scattering off the
  ions without being captured: the electrons are free before the interaction and
  remain free afterwards. In this process, the electron retains most of its
  pre-scattering energy. An electron, while passing by a much more massive singly
  ionized hydrogen atom, produces a radio photon of $\nu = 10$\,GHz. Calculate
  the mean electron thermal energy in the HII region, for the given temperature
  of the Str\"omgren sphere. Is this an example of free--free emission (answer
  YES or NO)? \hfill[4 points]
\item[(e)] Since the HII region is in local thermodynamic equilibrium, one can
  calculate the absorption coefficient that is proportional to the optical depth
  $\tau_\nu \propto \nu^{-2.1}$, and it turns out that at sufficiently high radio
  frequencies the hot plasma is nearly transparent and hence
  $\tau_\nu \ll 1$. The flux density of photons has a power-law spectrum of the
  form $S_\nu \propto \nu^{\beta}$. Find $\beta$ for the radio frequencies.
  \hfill[4 points]
\end{description}

\probhead{11.11}{Galactic Rotation}{2025}{Mumbai, India}{TH T06}{20}
% source id: 2025_TH_T06   tag: HI rotation curve, Oort constants
The rotation curve of our Galaxy is determined using line-of-sight velocity
measurements of neutral hydrogen (HI) clouds along various Galactic longitudes,
observed through the 21\,cm HI line.

Consider an HI cloud with Galactic longitude $l$, located at a distance $R$ from
the Galactic Centre (GC) and a distance $D$ from the Sun. Consider the Sun to be
at a distance $R_0 = 8.5$\,kpc from the GC. Assume that both the Sun and the HI
cloud are in circular orbits around the GC in the Galactic plane, with angular
velocities $\Omega_0$ and $\Omega$, and rotational velocities $V_0$ and $V$,
respectively.

The line-of-sight velocity ($V_r$) and transverse velocity ($V_t$) components of
the cloud, as observed from the Sun, can be expressed as
\begin{align*}
  V_r &= (\Omega - \Omega_0) R_0 \sin l \\
  V_t &= (\Omega - \Omega_0) R_0 \cos l - \Omega D
\end{align*}
Seen from the North Galactic Pole, the Galactic rotation is clockwise.
Throughout this problem, we shall take line-of-sight velocity to be positive when
receding and clouds will be treated as point objects.

\begin{description}
\item[(T06.1)] In the graph provided, sketch $V_r$ as a function of $D$ from
  $D = 0$ to $D = 2R_0$ for two lines of sight defined by (i) $l = 45^\circ$ and
  (ii) $l = 135^\circ$. Label each of your lines/curves with the value of $l$.
  \hfill[5]
\item[(T06.2)] The graph shows the average radial (solid, red curve) and
  transverse (dashed, blue curve) velocity components of stars at a distance of
  100\,pc from the Sun, plotted as a function of Galactic longitude. Using the
  graph, estimate the Sun's orbital period ($P$) around the GC in mega-years
  (Myr). \hfill[3]
\item[(T06.3)] Jan Oort noted that in the solar neighbourhood ($D \ll R_0$), the
  difference in angular velocities ($\Omega - \Omega_0$) will be small, and hence
  derived the following first order approximation for the line-of-sight and the
  transverse velocity components:
  \begin{align*}
    V_r &= A D \sin 2l \\
    V_t &= A D \cos 2l + B D
  \end{align*}
  where $A$ and $B$ are known as Oort's constants. Let us consider two cases:
  (I) the actual observed rotation curve of the Galaxy, and (II) the rotation
  curve for a hypothetical scenario where the Galaxy is devoid of dark matter and
  the whole mass of the Galaxy is assumed to be concentrated at its centre.
  \begin{description}
  \item[(T06.3a)] Derive expressions for the radial gradient of the rotational
    velocity at the location of the Sun, $\left.dV/dR\right|_{R=R_0}$, for the
    two cases. \hfill[2]
  \item[(T06.3b)] Express $A$ and $B$ in terms of $V_0$, $R_0$, and the radial
    gradient of rotational velocity at the location of the Sun,
    $\left.dV/dR\right|_{R=R_0}$. \hfill[8]
  \item[(T06.3c)] The ratios $(A/B)$ of Oort's constants for the two given
    cases, (I) and (II), are defined as $F_I$ and $F_{II}$ respectively.
    Determine $F_I$ and $F_{II}$. \hfill[2]
  \end{description}
\end{description}
\ioaafig{2025_TH_T06-f1.pdf}

\section*{Data-analysis problems}

\probhead{11.12}{}{2013}{Volos, Greece}{DA D4}{}
% source id: 2013_DA_D4   tag: Hyades distance moving-cluster method
Calculate the distance of the Hyades cluster using the moving cluster method.

You are given a list of 35 stars from the field of the Hyades open cluster,
observed by the Hipparcos space telescope. The information listed in the columns
of the data file for each of the 35 stars is: (a) the Hipparcos catalogue number
(HIP); (b) right ascension $\alpha$ [h m s]; (c) declination $\delta$
[$^\circ$ $'$ $''$]; (d) trigonometric parallax $\pi$ [$'' \times 10^3$];
(e) proper motion in right ascension multiplied by $\cos\delta$
($\mu_\alpha\cos\delta$) [$'' \times 10^3$/yr]; (f) proper motion in
declination ($\mu_\delta$) [$'' \times 10^3$/yr]; (g) radial velocity $v_r$
[km\,s$^{-1}$].
\ioaafig{2013_DA_D4-f1.pdf}

\probhead{11.13}{Black Hole in Milky Way}{2014}{Suceava, Romania}{DA D1}{}
% source id: 2014_DA_D1   tag: S-star orbit plot, Sgr A* mass
Using the provided observations of the orbits of S-stars around the compact radio
source Sgr\,A* at the centre of the Milky Way, determine the orbital elements of
the star and hence the mass of the central black hole.
\ioaafig{2014_DA_D1-f1.pdf}
\ioaafig{2014_DA_D1-f2.pdf}
\ioaafig{2014_DA_D1-f3.pdf}

\probhead{11.14}{Hypervelocity stars}{2020}{GeCAA (Estonia, remote)}{DA 3}{55}
% source id: 2020_DA_3   tag: HVS1 velocity, Galactic escape
Hypervelocity stars (HVSs) are stars moving fast enough to escape the
gravitational pull of the Galaxy. Using the provided data, determine the velocity
of the hypervelocity star and assess whether it is gravitationally bound to the
Milky Way.
\ioaafig{2020_DA_3-f1.pdf}
\ioaafig{2020_DA_3-f2.pdf}

\probhead{11.15}{Shapley Hypothesis}{2024}{Rio de Janeiro, Brazil}{DA D2}{75}
% source id: 2024_DA_D2   tag: globular clusters, distance to Galactic Centre
In 1918 Harlow Shapley used the spatial distribution of globular clusters to
argue that the Sun is not at the centre of the Milky Way. Using the provided
catalogue of globular clusters, reproduce this argument and estimate the distance
from the Sun to the Galactic Centre.
\ioaafig{2024_DA_D2-f1.pdf}

\section*{Observation --- telescope}

\probhead{11.16}{Identify open/globular clusters and nebulae}{2012}{Rio de Janeiro, Brazil}{OBS-TEL 4}{}
% source id: 2012_OBS-TEL_4   tag: visual object classification
Identify objects pointed out by the evaluator as an Open Cluster (OC), Globular
Cluster (GC), Emission Nebula (EN) or Planetary Nebula (PN), for each of four
objects.

\clearpage
\section*{Solutions to Chapter 11}

\solhead{11.1}{}
% source id: 2008_TH_S6
Rest frequency $\nu_o = 1420.41$\,MHz

Detected frequency $\nu = 1421.23$\,MHz
\[
  \Delta\nu = \nu_o - \nu = 1420.41 - 1421.23 = -0.82\ \mathrm{MHz}
\]
\hfill(25\%)

Using Doppler's shift, the speed of the cloud is:
\[
  v = \frac{\Delta\nu}{\nu_o}\,c
    = \frac{-0.82}{1421.23}\left(3\times10^{10}\right)
    = -1.73\times10^{7}\ \mathrm{cm/s}
\]
\hfill(15\%)

Since $v$ is negative, then the cloud is moving toward us. \hfill(10\%)

(Since $\nu > \nu_o$, then the cloud is moving toward us)

If $\mathcal{M}$ is mass of black hole, $v$ is the speed of the cloud and $R$
is the orbital radius of cloud, then
\[
  \mathcal{M} = \frac{Rv^2}{G}
\]
\hfill(25\%)

$R = 0.2\,\mathrm{pc} = 0.2 \times (3.086 \times 10^{18}\,\mathrm{cm})
= 6.17 \times 10^{17}\,\mathrm{cm}$, and from the Table of Astronomical and
Physical constants: $G = 6.67 \times 10^{-8}\,\mathrm{cm}^3/\mathrm{s}^2\,\mathrm{g}$

Then,
\[
  \mathcal{M} = \frac{Rv^2}{G}
  = \frac{\left(6.17\times10^{17}\right)\left(-1.73\times10^{7}\right)^2}{6.67\times10^{-8}}
  = 2.78\times10^{39}\ \mathrm{gr} = 1.40\times10^{6}\ \mathcal{M}_\odot
\]
\hfill(25\%)

\solhead{11.2}{Apparent number density of stars in the Galaxy}
% source id: 2009_TH_P17
\begin{description}
\item[a)] Relation between the apparent and absolute magnitude is given by
  \[
    m = M + 5\log\left(\frac{d}{10}\right)
  \]
  \hfill(3 points)

  where $d$ is in terms of parsec. Substituting $m = 18$ and $M = -0.2$,
  results in
  \[
    d = 4.37\times10^{4}\ \mathrm{pc}
  \]
  \hfill(5 points)
\item[b)] Adding the term for the extinction, changes the magnitude distance
  relation as follows
  \[
    m = M + 0.7x + 5\log\,(100x)
  \]
  where $x$ is given in terms of kilo parsec. To have a rough value for $x$,
  after substituting $m$ and $M$, this equation reduces to
  \[
    8.2 = 0.7x + 5\log\,(x)
  \]
  \hfill(6 points)

  To solve this equation, we examine
  \[
    x = 5\,,\,5.5\,,\,6\,,\,6.5
  \]
  where the best value is obtained roughly $x \cong 6.1\ \mathrm{kpc}$.
  \hfill(8 points)
\item[c)] For a solid angle $\Omega$, the number of observed red clump stars
  at the distance in the range of $x$ and $x + \Delta x$ is given by
  \[
    \Delta N = \Omega\,x^2\,n(x)f\,\Delta x
  \]
  So the number of stars observed in $\Delta x$ is given by
  \[
    \frac{\Delta N}{\Delta x} = \Omega\,x^2\,n(x)f
  \]
  \hfill(6 points)

  From the relation between the distance and apparent magnitude we have
  \begin{align*}
    m_1 &= M + 5\log\left(\frac{x}{10}\right)\\
    m_2 &= M + 5\log\left(\frac{x+\Delta x}{10}\right)\\
    \Delta m &= 5\log\left(\frac{x+\Delta x}{x}\right)\\
    \Delta m &= 5\log\left(1 + \frac{\Delta x}{x}\right)\\
    \Delta m &= \frac{5}{\ln 10}\,\ln\left(1 + \frac{\Delta x}{x}\right)\\
    \Delta m &= \frac{5}{\ln 10}\left(\frac{\Delta x}{x}\right)
  \end{align*}
  Replacing $\Delta x$ with $\Delta m$, results in
  \[
    \frac{\Delta N}{\Delta m}
      = \frac{\Delta N}{\Delta x}\times\frac{\Delta x}{\Delta m}
  \]
  So the number of stars for a given magnitude is obtained by \hfill(5 points)
  \[
    \frac{\Delta N}{\Delta m} = \frac{\Omega \ln 10}{5}\,n(x)x^3 f
  \]
  Finally we substituting % sic
  $x$ in terms of apparent magnitude using $x = 10^{\frac{m+5.2}{5}}$.
  % sic: the exponent is printed with m+5.2 here but m-5.2 in the final
  % expressions below; with x in kpc the correct form is 10^{(m-5.2)/5}.
  In the case of no extinction, we are able to observe the Galaxy beyond the
  center. So $\frac{dN}{dm}$ has two terms in $x < R_0$ and $x > R_0$. The
  relation between $x$ and $r$ for these two cases are
  \[
    x = R_0 - r \qquad x < R_0
  \]
  and
  \[
    x = R_0 + r \qquad x > R_0
  \]
  \hfill(6 points)

  So in general we can write $\frac{\Delta N}{\Delta m}$ as
  \[
    \frac{\Delta N}{\Delta m}
      = \frac{\Omega \ln 10}{5}\,n_0
        \exp\Big(\frac{10^{\frac{m-5.2}{5}}}{R_d}\Big)
        \times 10^{\frac{3(m-5.2)}{5}} f
    \qquad x < R_0
  \]
  \[
    \frac{\Delta N}{\Delta m}
      = \frac{\Omega \ln 10}{5}\,n_0
        \exp\Big(\frac{2R_0}{R_d}\Big)
        \exp\Big(-\frac{10^{\frac{m-5.2}{5}}}{R_d}\Big)
        \times 10^{\frac{3(m-5.2)}{5}} f\,\Theta(x_0 - x)
    \qquad x > R_0
  \]
  \hfill(6 points)

  where $\Theta(x)$ is the step function and $x_0$ is the maximum observable
  distance.
\end{description}

\solhead{11.3}{}
% source id: 2009_TH_P8
\[
  \frac{1}{2}Mv_{esc}^2 = \frac{1}{2}M\,\frac{2GM}{R}
\]
\hfill(4 points)

The velocity must be smaller than the escape velocity
($v_{esc} \approx \sqrt{2}\,v_{rms}$) and since the problem is an estimation,
any velocities smaller than escape velocity is accepted and therefore the
escape velocity can be replaced by $v_{rms}$.
\[
  Mv_{rms}^2 = \frac{GM^2}{R}
\]
\hfill(2 points)
\[
  M = \frac{Rv_{rms}^2}{G}
    = \frac{20 \times 3.09\times10^{16} \times 9\times10^{6}}{6.67\times10^{-11}}
    = 8.3\times10^{34}\ \mathrm{Kg}
\]
\[
  M = 4.2\times10^{4}\,M_\odot
\]
\hfill(4 points)

\solhead{11.4}{}
% source id: 2009_TH_P9
\ioaafig{2009_TH_P9-s1.pdf}
In the figure, S is the Sun and $R_0$ and $V_0$ are Sun distance and velocity.
The distance and velocity of star $P$ is denoted by $R$ and $V = V_0$. The
radial velocity of star $P$ respect to the Sun is
\[
  V_r = V\cos\alpha - V_0\sin l = V_0(\cos\alpha - \sin l)
\]
\hfill(4 points)

In SCP triangle we have
\[
  \frac{\sin l}{R} = \frac{\cos\alpha}{R_0}
  \;\Rightarrow\; \cos\alpha = \frac{R_0}{R}\sin l
\]
\hfill(1 point)

So
\[
  V_r = V_0\left(\frac{R_0}{R} - 1\right)\sin l
\]
\hfill(2 points)
\begin{align*}
  \frac{R_0}{R} - 1 &= \frac{V_r}{V_0\sin l}\\
  \frac{R_0}{R} &= \frac{V_r + V_0\sin l}{V_0\sin l}\\
  R &= R_0\,\frac{V_0\sin l}{V_r + V_0\sin l}
\end{align*}
\hfill(2 points)
\begin{align*}
  R &= 8\times10^{3}\,\frac{250\sin 15^\circ}{100 + 250\sin 15^\circ}\\
  R &= 3\times10^{3}\ \mathrm{pc}
\end{align*}
\hfill(1 point)

\solhead{11.5}{}
% source id: 2010_TH_S10
Since $h_z = 300\,\mathrm{pc}$, we can substitute this into the vertical
(exponential) disc equation:
\[
  n(0.5\,\mathrm{kpc}) = n_0 \exp(-\left|500\,\mathrm{pc}\right|/300\,\mathrm{pc})
  \approx 0.189\,n_0
\]
\hfill(5 points)

In other words, the density of G-type MS stars at 0.5\,kpc above the plane is
just under 19\% of its mid-plane value.

For $z = 1.5\,\mathrm{kpc}$, this works out as $0.007$. \hfill(5 points)

\solhead{11.6}{}
% source id: 2011_TH_S2
$v_e = \sqrt{2GM_{\mathrm{sol}}\,n/R}$, \hfill[5]

solve for $n$. ($\sim 100\,000$). \hfill[5]

\solhead{11.7}{}
% source id: 2013_TH_S10
\[
  F = -\frac{GM_{BH}M_*}{R^2} = -\frac{M_*v^2}{R}
\]
\hfill(2 Point)

But $v = \dfrac{2\pi R}{P}$. Therefore
\[
  \frac{GM_{BH}}{R^2} = 4\pi^2\,\frac{R}{P^2}
  \qquad\text{or}\qquad
  GM_{BH} = 4\pi^2\,\frac{R^3}{P^2}
\]
\hfill(2 Points)

Similarly:
\[
  GM_\odot = 4\pi^2\,\frac{(1\,AU)^3}{(1\,yr)^2}
  \qquad\text{or}\qquad
  G = \frac{1}{M_\odot}\,4\pi^2\,\frac{(1\,AU)^3}{(1\,yr)^2}
\]
\hfill(2 Point)

From Kepler's 3rd law we get:
\[
  \frac{M_{BH}}{M_\odot} = \frac{(R/1\,AU)^3}{(P/1\,yr)^2}
\]
\hfill(1 Point)

Inserting the given data we find the distance of the star from the black hole:
\[
  R = \frac{0.12}{200.000}\,(8000)(3\times10^{18}\,\mathrm{cm})
    = 1.4\times10^{16}\,\mathrm{cm} = 960\,AU
\]
\hfill(2 Point)

Therefore:
\[
  \frac{M_{BH}}{M_\odot} = \frac{(960)^3}{(15)^2} = 4\times10^{6}
\]
form which % sic
we calculate the mass of the black hole:
\[
  M_{BH} = 4\times10^{6}\,M_\odot
\]
\hfill(1 Point)

\solhead{11.8}{}
% source id: 2015_TH_S10
The molecules rotate with angular velocity $\omega$ and for spherical
molecules, their moment of inertia is
\begin{align*}
  I_H &= I_{cm} + m_H r_H^2\\
      &= \frac{2}{5}m_H^2 + m_H r_H^2 = \frac{7}{5}m_H r_H^2
      % sic: the first term is printed as (2/5)m_H^2; it should read
      % (2/5)m_H r_H^2
\end{align*}
\hfill(30 points)

For $H_2$ molecule, yield
\[
  I_{H_2} = 2I_H = \frac{14}{5}m_H r_H^2\,,
\]
Thus
\[
  I = 6.401\times10^{-48}\ \mathrm{Kg\,m^2}
\]

In thermal equilibrium (only 2 rotational degrees of freedom), we have:
\[
  \frac{1}{2}I\omega^2 = kT
\]
\hfill(30 points)

So we can estimate:
\[
  \omega = \sqrt{2kT/I}
  = (2 \times 1.38\times10^{-16} \times 115/6.40\times10^{-48})^{1/2}
  = 2.23\times10^{13}\,s^{-1}
  % sic: the source prints 1.38e-16 (CGS k) although the numerical result
  \]
\[
  \nu = \frac{\omega}{2\pi} = 3.55\times10^{13}\,\mathrm{Hz}
\]
\hfill(40 points)

\solhead{11.9}{The Milky Way's Distant Outer Arm}
% source id: 2017_TH_T3
\begin{description}
\item[a)]
  \ioaafig[0.54]{2017_TH_T3-s1.pdf}
  If there is the flat rotation curve beyond 2\,kpc from the Galactic centre,
  the orbital velocity of the Sun and the start of the arm are both equal to
  \[
    v_r = -v_{\mathrm{LSR}}\sin(\ell) + v_{\mathrm{LSR}}\sin(\alpha)
  \]
  \hfill[3.0]

  Sine law ($\triangle AGS$):\quad
  $\dfrac{\sin(\ell)}{d} = \dfrac{\sin(\alpha)}{R_0}$ \hfill[1.0]
  \[
    v_r = -v_{\mathrm{LSR}}\sin(\ell) + v_{\mathrm{LSR}}R_0\,\frac{\sin(\ell)}{d}
  \]
  \hfill[1.0]
  \[
    20.9 = 220\sin(13.25^\circ) - 220\times8.5\times\frac{\sin(13.25^\circ)}{d}
  \]
  \hfill[1.0]

  \emph{Correctly substitute numerical values}
  \[
    d = 14.5\,\mathrm{kpc}
  \]
  \hfill[1.0]
\item[b)]
  Cosine law ($\triangle AGS$):
  \[
    d^2 = x^2 + R_0^2 - 2xR_0\cos(\ell)
  \]
  \hfill[1.0]
  \[
    14.5^2 = x^2 + 8.5^2 - 2\times8.5\,x\cos(13.25^\circ)
  \]
  \hfill[1.0]

  \emph{Correctly substitute numerical values}
  \[
    x = \frac{17\cos(13.25^\circ)
        + \sqrt{17^2\cos^2(13.25^\circ) - 4\times(8.5^2 - 14.5^2)}}{2}
  \]
  \[
    x = 22.6\,\mathrm{kpc}
  \]
  \hfill[1.0]
\end{description}

\solhead{11.10}{HII region}
% source id: 2020_TH_5
\begin{description}
\item[(a)] As the number of H-atoms undergoing ionization and recomination % sic
  are balanced at $R_S$, each photon can ionize exactly one hydrogen atom and
  each neutral hydrogen has exactly one proton and one electron,
  \begin{align*}
    n_{\mathrm{recomb}} &= n_{HII} = Q\\
    \text{and } n_e &= n_p = n_H \tag*{(2.0)}\\
    n_{\mathrm{recomb}} &= \alpha n_p n_e V_S\\
    Q &= \alpha n_H^2\,\frac{4\pi}{3}R_S^3 \tag*{(1.0)}\\
    \therefore R_S &= \sqrt[3]{\frac{3Q}{4\pi\alpha n_H^2}} \tag*{(1.0)}\\
      &= \sqrt[3]{\frac{3\times10^{49}}{4\pi\times10^{-19}\times(10^{8})^2}}\\
      &= 1.3\times10^{17}\,\mathrm{m}\\
    \therefore R_S &\approx 4\,\mathrm{pc} \tag*{(1.0)}
  \end{align*}
\item[(b)] Per one unit length distance, a typical photon will encounter
  $\sigma n_H$ H-atoms. Thus, the mean free path will be,
  \begin{align*}
    l_{\nu_0} &= \frac{1}{\sigma n_H} = \frac{1}{\sigma 10^{-21}\times10^{8}}
    % sic: the source prints "\sigma 10^{-21}" in the denominator
    \\
    l_{\nu_0} &= 10^{13}\,\mathrm{m} \tag*{(2.0)}\\
    \therefore l_{\nu_0} &\approx 10^{-4}R_S \ll R_S
  \end{align*}
  Thus, the boundary layer of the sphere is very thin as compared to its total
  size. Hence, ``YES'' this ionized nebula is very sharp-edged. \hfill(1.0)
\item[(c)] For Stromgren sphere to form, all H-atoms ($N$) inside $R_S$ need
  to be ionized. Thus, time $t_S$ required will be
  \begin{align*}
    N &= V_S n_H = \frac{4\pi}{3}R_S^3 n_H \tag*{(1.0)}\\
    t_S &= \frac{N}{Q} = \frac{4\pi R_S^3 n_H}{3Q} \tag*{(1.0)}\\
        &= \frac{1}{\alpha n_H} = \frac{1}{10^{-19}\times10^{8}}\\
        &= 10^{11}\,\mathrm{s}\\
        &\approx 3000\,\mathrm{yr} \tag*{(2.0)}
  \end{align*}
  For typical nebular densities, the main-sequence lifetime of LBV stars is
  much longer than this ionization time, and hence our assumption of a
  stationary system is justified.
\item[(d)] The mean electron thermal energy in a plasma of temperature
  $T_e = 10^4$\,K is
  \[
    E_e \approx k_B T_e \approx 1.4\times10^{-19}\,\mathrm{J}
        \approx 0.9\,\mathrm{eV}
  \]
  \hfill(2.0)

  The energy of a photon is
  \[
    E_\gamma = h\nu \approx 6.6\times10^{-24}\,\mathrm{J}
             \approx 4\times10^{-5}\,\mathrm{eV}
  \]
  \hfill(1.0)

  As the energy of the radio photon is much much smaller than that of the
  electron, the answer is ``Yes''. \hfill(1.0)
\item[(e)] As the plasma is nearly transparent, using the Rayleigh--Jeans law,
  \begin{align*}
    B_\nu &= \frac{2k_B T_e \nu^2}{c^2} \tag*{(1.0)}\\
    S_\nu &\propto B_\nu \tau_\nu = \frac{2k_B T_e \nu^2}{c^2}\,\tau_\nu\\
    S_\nu &\propto \nu^2\cdot\nu^{-2.1}\\
    \therefore S_\nu &\propto \nu^{-0.1} \tag*{(3.0)}
  \end{align*}
\end{description}

\solhead{11.11}{Galactic Rotation}
% source id: 2025_TH_T06
\begin{description}
\item[(T06.1)]
  \ioaafig[0.66]{2025_TH_T06-s1.pdf}
  The concept of differential rotation of the Galaxy is used here.
  \begin{itemize}
  \item For $l = 45^\circ$:
    \begin{itemize}
    \item From Sun to the GC, the quantity $(\Omega - \Omega_0)$ increases,
      which implies $V_{\mathrm{r}}$ increases.
    \item As shown in the figure, the line of sight intersects circular orbits
      that are at varying distances from the GC.
    \item The line of sight is closest to the GC when
      $D = R_0\cos l = 6.01$\,kpc. $V_{\mathrm{r}}$ reaches the maximum value
      here.
    \item Beyond this, $(\Omega - \Omega_0)$ decreases implying
      $V_{\mathrm{r}}$ decreases.
    \item $V_{\mathrm{r}} = 0$ when the line of sight intersects the Sun's
      orbit. Here, $D = 2R_0\cos l = 12.02$\,kpc
    \item Beyond the Sun's orbit, $\Omega < \Omega_0$. $V_{\mathrm{r}}$ is
      negative and its magnitude increases with increasing $D$.
    \end{itemize}
  \item $l = 135^\circ$
    \begin{itemize}
    \item The line of sight lies outside the Sun's orbit. Here,
      $V_{\mathrm{r}}$ is negative, there is no peak, and its magnitude
      increases with increasing $D$.
    \end{itemize}
  \end{itemize}
  The plot of the line-of-sight velocity as a function of the distance from
  the Sun for two values of $l$ is as follows:
  \ioaafig{2025_TH_T06-s2.pdf}
\item[(T06.2)] Consider $l = 90$,
  \begin{align*}
    V_{\mathrm{r}} &= 0 \text{ implying } \Omega = \Omega_0
       \tag*{(0.5)}\\
    V_{\mathrm{t}} &= -\Omega_0 D = -2.75\,\mathrm{km\,s^{-1}}
       \tag*{(0.5)}\\
    \Omega_0 &= 8.91\times10^{-16}\,\mathrm{rad\,s^{-1}}
       \tag*{(1.0)}\\
    P &= \frac{2\pi}{\Omega_0} = 224\,\mathrm{Myr}
       \tag*{(1.0)}
  \end{align*}
\item[(T06.3a)] \mbox{}
  \begin{description}
  \item[(I)] The observed rotation curve is flat near the Sun.

    $V$ is constant $\rightarrow$
    $\left.\dfrac{dV}{dR}\right|_{R=R_0} = 0$. \hfill(1.0)
  \item[(II)] If the whole mass of the Galaxy is assumed to be concentrated
    at the centre, then we consider Keplerian motion.
    \[
      V = \sqrt{\frac{GM}{R}}
      \;\rightarrow\;
      \left.\frac{dV}{dR}\right|_{R=R_0} = -\frac{1}{2}\frac{V_0}{R_0}.
    \]
    \hfill(1.0)
  \end{description}
\item[(T06.3b)] In the solar neighbourhood, $(\Omega - \Omega_0)$ will be
  very small. Using the approximation
  \[
    f(x) \simeq f(x_0)
      + \left.\frac{df}{dx}\right|_{x=x_0}(x - x_0),
    \text{ for } x \approx x_0, \text{ given in the Data Sheet,}
  \]
  \begin{align*}
    (\Omega - \Omega_0)
      &= \left.\frac{d\Omega}{dR}\right|_{R=R_0}(R - R_0)
       \tag*{(2.0)}\\
      &= \left[\left(\frac{1}{R}\frac{dV}{dR}\right)\bigg|_{R=R_0}
         - \left(\frac{1}{R^2}V\right)\bigg|_{R=R_0}\right](R - R_0)\\
      &= \frac{1}{R_0^2}\left[R_0\left.\frac{dV}{dR}\right|_{R=R_0}
         - V_0\right](R - R_0)
  \end{align*}
  Thus the line-of-sight velocity expression becomes
  \[
    V_{\mathrm{r}} = \frac{1}{R_0^2}
      \left[R_0\left.\frac{dV}{dR}\right|_{R=R_0} - V_0\right]
      (R - R_0)R_0\sin l
  \]
  Further, in the solar neighbourhood, $D \ll R_0$. So we can approximate
  \[
    (R - R_0) \approx -D\cos l
  \]
  \hfill(1.0)

  Thus, the line-of-sight velocity can be expressed as,
  \begin{align*}
    V_{\mathrm{r}}
      &= \frac{1}{R_0^2}
         \left[R_0\left.\frac{dV}{dR}\right|_{R=R_0} - V_0\right]
         (-D\cos l)R_0\sin l\\
      &= \frac{1}{2}
         \left[\frac{V_0}{R_0} - \left.\frac{dV}{dR}\right|_{R=R_0}\right]
         D\sin 2l
  \end{align*}
  Comparing with $V_{\mathrm{r}} = AD\sin 2l$, we get
  \[
    A = \frac{1}{2}
        \left[\frac{V_0}{R_0} - \left.\frac{dV}{dR}\right|_{R=R_0}\right]
  \]
  \hfill(2.0)

  Also, in the solar neighbourhood, $\Omega D \approx \Omega_0 D$. Using this
  approximation and $(R - R_0) \approx -D\cos l$, the tangential velocity is
  \hfill(1.0)
  \begin{align*}
    V_{\mathrm{t}} &= (\Omega - \Omega_0)R_0\cos l - \Omega D\\
      &= \frac{1}{R_0^2}
         \left[R_0\left.\frac{dV}{dR}\right|_{R=R_0} - V_0\right]
         (-D\cos l)R_0\cos l - \Omega_0 D\\
      &= \left[\frac{V_0}{R_0}
         - \left.\frac{dV}{dR}\right|_{R=R_0}\right]D\cos^2 l - \Omega_0 D\\
      &= \frac{1}{2}\left[\frac{V_0}{R_0}
         - \left.\frac{dV}{dR}\right|_{R=R_0}\right]D(1 + \cos 2l)
         - \frac{V_0}{R_0}D\\
      &= AD\cos 2l
         + \frac{1}{2}\left[\frac{V_0}{R_0}
         - \left.\frac{dV}{dR}\right|_{R=R_0}\right]D - \frac{V_0}{R_0}D\\
      &= AD\cos 2l
         - \frac{1}{2}\left[\frac{V_0}{R_0}
         + \left.\frac{dV}{dR}\right|_{R=R_0}\right]D
  \end{align*}
  Comparing with $V_{\mathrm{t}} = AD\cos 2l + BD$, we get
  \[
    B = -\frac{1}{2}
        \left[\frac{V_0}{R_0} + \left.\frac{dV}{dR}\right|_{R=R_0}\right]
  \]
  \hfill(2.0)
\item[(T06.3c)] For a flat rotation curve
  $\left(\dfrac{dV}{dR} = 0\right)$, $F_{\mathrm{I}} = -1$. \hfill(1.0)

  For Keplerian rotation
  $\left(\dfrac{dV}{dR} = -\dfrac{1}{2}\dfrac{V}{R}\right)$,
  $F_{\mathrm{II}} = -3$. \hfill(1.0)
\end{description}

\solhead{11.12}{}
% source id: 2013_DA_D4
\begin{description}
\item[(1)] The student should be able to import the ascii data of the
  \emph{txt} file in the \emph{MS Excel} spreadsheet application.
  \hfill(12 Points)
\item[(2)] In order to calculate the angular distance $\varphi$ (next step),
  all coordinates should be converted into decimal degrees. The students
  should be able to convert the [hours, min, sec] and the
  [$^\circ$,\ $'$,\ $''$] into decimal degrees, using. % sic
  For the right ascension:
  $\alpha_{\deg} = h\times15 + \dfrac{m}{60} + \dfrac{s}{3600}$
  and for the declination:
  $\delta_{\deg} = {}^\circ + \dfrac{'}{60} + \dfrac{''}{3600}$.
  All coordinates are positive, so the student does not have to worry about
  checking the sign, which is not trivial. \hfill(12 Points)
\item[(3)] Using the cosine law,
  $\varphi = \arccos(\sin\delta_1\times\sin\delta_2
    + \cos\delta_1\times\cos\delta_2\times\cos\,(a_1 - a_2))$,
  the angular distance, $\varphi$, between each of the stars and the point of
  convergence should be calculated \hfill(15 Points)
\item[(4)] The student should easily calculate the proper motion of each
  star, by inserting the given equation,
  $\mu = \sqrt{(\mu_a\cos\delta)^2 + \mu_d^2}$, in the spreadsheet.
  \hfill(9 Points)
\item[(5)] Again, inserting the given equation,
  $r_\mu = \dfrac{v_r\tan\varphi}{4.74047\,\mu}$ in the spreadsheet
  (remembering to divide the given values of $\mu$ by 1000 to get arcseconds),
  the student should be able to calculate the distance, $r_\mu$, of each star.
  Any star whose distance from the centre of the cluster is larger than
  10\,pc should be omitted from the following calculations. \hfill(12 Points)
\item[(6)] The trigonometric parallax distance is given by
  $r_\pi = \dfrac{1}{\pi\times10^{-3}}$. This equation should be inserted in
  the spreadsheet. The result is in parsecs [pc]. \hfill(6 Points)
\item[(7)] Using the statistical functions \emph{AVERAGE} and \emph{STDEV} of
  MS Excel, the student should easily calculate the \emph{average} and the
  \emph{standard deviation} of $r_\mu$ and $r_\pi$. \hfill(6 Points)
\item[(8)] The method with the smaller standard deviation is, obviously more
  accurate. \hfill(3 Point)
\end{description}

(\emph{macros} are allowed)

\solhead{11.13}{Black Hole in Milky Way}
% source id: 2014_DA_D1
\textbf{Detailed solution}

Usually the plane of apparent orbit (sky plane) is different from the plane of
the relative orbit due to the angle between the two planes ($i = 0$). % sic:
% the text prints (i = 0) here although it describes the general case of a
% nonzero inclination.
\ioaafig{2014_DA_D1-s1.pdf}

In the case of the problem the sky plane and the plane of the relative orbit
are the same ($i = 0$). See the figure.
\ioaafig{2014_DA_D1-s2.pdf}

The linear position coordinates $(u, w)$, of the projection of the star S* in
the sky plane with the coordinate axes (U, W), having the origin in
Sagittarius A*, expressed in light-days are given in the table 1. (for an
angular distance $\varphi = 1\,\mathrm{arcsec}$ correspond in sky plane the
linear distance $d = 41$\,light day).

\begin{center}
Table 1 The coordinates in light-days

\begin{tabular}{|c|c|c|c|c|}
\hline
 & Date (year) & $\alpha$(arcsec) & $u$(ld) & $\beta$(arcsec) \quad $w$(ld)\\
\hline
1  & 1995.222 & 0.117  & 4.797  & $-0.166$ \quad $-6.806$\\
2  & 1997.526 & 0.097  & 3.977  & $-0.189$ \quad $-7.749$\\
3  & 1998.326 & 0.087  & 3.567  & $-0.192$ \quad $-7.872$\\
4  & 1999.041 & 0.077  & 3.157  & $-0.193$ \quad $-7.913$\\
5  & 2000.414 & 0.052  & 2.132  & $-0.183$ \quad $-7.503$\\
6  & 2001.169 & 0.036  & 1.476  & $-0.167$ \quad $-6.847$\\
7  & 2002.831 & $-0.000$ & 0.000  & $-0.120$ \quad $-4.920$\\ % sic: comma
8  & 2003.584 & $-0.016$ & $-0.656$ & $-0.083$ \quad $-3.403$\\
9  & 2004.165 & $-0.026$ & $-1.066$ & $-0.041$ \quad $-1.681$\\
10 & 2004.585 & $-0.017$ & $-0.697$ & 0.008 \quad 0.328\\
11 & 2004.655 & $-0.004$ & $-0.164$ & 0.014 \quad 0.574\\
12 & 2004.734 & 0.008  & 0.328  & 0.017 \quad 0.697\\
13 & 2004.839 & 0.021  & 0.861  & 0.012 \quad 0.492\\
14 & 2004.936 & 0.037  & 1.517  & 0.009 \quad 0.369\\
15 & 2005.503 & 0.072  & 2.952  & $-0.024$ \quad $-0.984$\\
16 & 2006.041 & 0.088  & 3.608  & $-0.050$ \quad $-2.050$\\
17 & 2007.060 & 0.108  & 4.428  & $-0.091$ \quad $-3.731$\\
\hline
\end{tabular}
\end{center}

To represent the image of the relative orbit of star S* in a plane parallel
with the sky plane we use a scale of
$S = 1200\,\dfrac{\mathrm{mm}}{\mathrm{arcsec}}$, and for conversion we use
$x = \alpha S$ and $y = \beta S$, see figure 3. These coordinates of the image
of the star S* relative to the origin of the axes -- the image of Sagittarius
A* are given in the below table 3
% sic: the tables are captioned "Table 2 The coordinates in scaled plane"
\ioaafig{2014_DA_D1-s3.pdf}

\begin{center}
Table 2 The coordinates in scaled plane

\begin{tabular}{|c|c|c|c|c|c|}
\hline
 & Date (year) & $\alpha$(arcsec) & $x$(mm) & $\beta$(arcsec) & $y$(mm)\\
\hline
1  & 1995.222 & 0.117  & 140.4  & $-0.166$ & $-199.2$\\
2  & 1997.526 & 0.097  & 116.4  & $-0.189$ & $-226.8$\\
3  & 1998.326 & 0.087  & 104.4  & $-0.192$ & $-230.4$\\
4  & 1999.041 & 0.077  & 92.4   & $-0.193$ & $-231.6$\\
5  & 2000.414 & 0.052  & 62.4   & $-0.183$ & $-219.6$\\
6  & 2001.169 & 0.036  & 43.2   & $-0.167$ & $-200.4$\\
7  & 2002.831 & $-0.000$ & 0.0    & $-0.120$ & $-144.0$\\
8  & 2003.584 & $-0.016$ & $-20.0$  & $-0.083$ & $-100.0$\\
9  & 2004.165 & $-0.026$ & $-32.0$  & $-0.041$ & $-50.0$\\
10 & 2004.585 & $-0.017$ & $-20.4$  & 0.008  & 9.6\\
11 & 2004.655 & $-0.004$ & $-4.8$   & 0.014  & 16.8\\
12 & 2004.734 & 0.008  & 9.6    & 0.017  & 20.4\\
13 & 2004.839 & 0.021  & 25.2   & 0.012  & 14.4\\
14 & 2004.936 & 0.037  & 44.4   & 0.009  & 10.8\\
15 & 2005.503 & 0.072  & 86.4   & $-0.024$ & $-28.8$\\
16 & 2006.041 & 0.088  & 106.0  & $-0.050$ & $-60.0$\\
17 & 2007.060 & 0.108  & 130.0  & $-0.091$ & $-110.0$\\
\hline
\end{tabular}
\end{center}

The following pictures represents the plot of the values in the table 2, and
the measurements which have to be done in order to find out the required data.
\textit{[orbit plot is in a separate millimetric-paper annex ('See Anex', 2014\_da\_sol.pdf p.21) not included in the source PDF]}

\textbf{Numeric processing}

\begin{description}
\item[a)] From the measurements done on paper system (X, Y):
  \[
    a_{maesured} \approx 135\,\mathrm{mm};\;
    % sic: "maesured"/"mesured" as printed
    b_{mesured} \approx 66\,\mathrm{mm};\;
    d_{\min} = 17\,\mathrm{mm},
  \]
  Corresponding to the numeric scale:
  \[
    a_{mesured} = \varphi_a\cdot S;\;
    b_{mesured} = \varphi_b\cdot S;\;
    d_{\min\,mes} = \varphi_{\min}\cdot S,
  \]
  \begin{align*}
    \varphi_a &= 0.1125\,\mathrm{arcsec};\\
    \varphi_b &= 0.055\,\mathrm{arcsec};\\
    \varphi_{d\min} &= 0.01416\,\mathrm{arcsec};
  \end{align*}
  \begin{align*}
    1\,\mathrm{arcsec}\ &\ \dots\dots\dots\dots\ 41\,\mathrm{light\ day};\\
    \varphi_a = 0.1125\,\mathrm{arcsec}\ &\ \dots\dots\dots\dots\ a;\\
    \varphi_b = 0.055\,\mathrm{arcsec}\ &\ \dots\dots\dots\dots\ b;\\
    \varphi_{d\min} = 0.014666\,\mathrm{arcsec}\ &\ \dots\dots\dots\dots\ r_{\min};
    % sic: 0,01416 above vs 0,014666 here, as printed
  \end{align*}
  \begin{align*}
    a &= \frac{0.1125\,\mathrm{arcsec}\cdot41\,\mathrm{light\ days}}
              {1\,\mathrm{arcsec}};\;
    a = 4.6125\,\mathrm{light\ days},\\
    b &= \frac{0.055\,\mathrm{arcsec}\cdot41\,\mathrm{light\ days}}
              {1\,\mathrm{arcsec}};\;
    b = 2.255\,\mathrm{light\ days},\\
    r_{\min} &= \frac{0.0146\,\mathrm{arcsec}\cdot41\,\mathrm{light\ days}}
              {1\,\mathrm{arcsec}} = 0.5986\,\mathrm{light\ days};
  \end{align*}
  \[
    c = \sqrt{a^2 - b^2} \approx 4.0236\,\mathrm{light\ days};
  \]
  \[
    r_{\min} = a - c = 0.5889\,\mathrm{light\ days};\quad
    r_{\max} = a + c = 8.6361\,\mathrm{light\ days};
  \]
  \[
    e = \sqrt{1 - \frac{b^2}{a^2}} \approx 0.87,
  \]
\item[b)] $S_{\mathrm{elipsa}} = 28560\,\mathrm{mm}^2$; See Anex % sic
  \begin{align*}
    T\ &\ \dots\dots\dots\dots\ S_{\mathrm{elipsa}};\\
    \Delta t\ &\ \dots\dots\dots\dots\ \Delta S;\\
    T &= \frac{S_{\mathrm{elipsa}}}{\Delta S}\cdot\Delta t;
  \end{align*}
  Displacement $1 \rightarrow 2$
  \begin{align*}
    \Delta S_{12} &= \frac{252\,\mathrm{mm}\times34\,\mathrm{mm}}{2}
      + 86\,\mathrm{mm}^2 \approx 4370\,\mathrm{mm}^2;\\
    t_1 &= 1995.222\,\mathrm{years};\; t_2 = 1997.526\,\mathrm{years};\\
    \Delta t_{12} &= t_2 - t_1 = 2.304\,\mathrm{years};\\
    T\ &\ \dots\dots\dots\dots\ 28560\,\mathrm{mm}^2;\\
    2.304\,\mathrm{years}\ &\ \dots\dots\dots\dots\ 4370\,\mathrm{mm}^2;\\
    T &= 15.057\,\mathrm{years};
  \end{align*}
  Displacement $2 \rightarrow 3$
  \begin{align*}
    \Delta S_{23} &= \frac{252\,\mathrm{mm}\times12\,\mathrm{mm}}{2}
      + 20\,\mathrm{mm}^2 \approx 1532\,\mathrm{mm}^2;\\
    t_2 &= 1997.526\,\mathrm{years};\; t_3 = 1998.326\,\mathrm{years};\\
    \Delta t_{23} &= t_3 - t_2 = 0.800\,\mathrm{years};\\
    T\ &\ \dots\dots\dots\dots\ 28560\,\mathrm{mm}^2;\\
    0.800\,\mathrm{years}\ &\ \dots\dots\dots\dots\ 1532\,\mathrm{mm}^2;\\
    T &= 14.913\,\mathrm{years};
  \end{align*}
  Displacement $3 \rightarrow 4$
  \begin{align*}
    \Delta S_{34} &= \frac{250\,\mathrm{mm}\times11\,\mathrm{mm}}{2}
      + 26\,\mathrm{mm}^2 \approx 1400\,\mathrm{mm}^2;\\
    t_3 &= 1998.326\,\mathrm{years};\; t_4 = 1999.041\,\mathrm{years};\\
    \Delta t_{34} &= t_4 - t_3 = 0.715\,\mathrm{years};\\
    T\ &\ \dots\dots\dots\dots\ 28560\,\mathrm{mm}^2;\\
    0.715\,\mathrm{years}\ &\ \dots\dots\dots\dots\ 1400\,\mathrm{mm}^2;\\
    T &= 14.586\,\mathrm{years};
  \end{align*}
  Displacement $4 \rightarrow 5$
  \begin{align*}
    \Delta S_{45} &= \frac{244\,\mathrm{mm}\times22\,\mathrm{mm}}{2}
      + 30\,\mathrm{mm}^2 = 2714\,\mathrm{mm}^2;\\
    t_4 &= 1999.041\,\mathrm{years};\; t_5 = 2000.414\,\mathrm{years};\\
    \Delta t_{45} &= t_5 - t_4 = 1.373\,\mathrm{years};\\
    T\ &\ \dots\dots\dots\dots\ 28560\,\mathrm{mm}^2;\\
    1.373\,\mathrm{years}\ &\ \dots\dots\dots\dots\ 2714\,\mathrm{mm}^2;\\
    T &= 14.448\,\mathrm{years};
  \end{align*}
  Displacement $5 \rightarrow 6$
  \begin{align*}
    \Delta S_{56} &= \frac{220\,\mathrm{mm}\times13\,\mathrm{mm}}{2}
      + 5\,\mathrm{mm}^2 = 1435\,\mathrm{mm}^2;\\
    t_5 &= 2000.414\,\mathrm{years};\; t_6 = 2001.169\,\mathrm{years};\\
    \Delta t_{56} &= 0.755\,\mathrm{years};\\
    T\ &\ \dots\dots\dots\dots\ 28560\,\mathrm{mm}^2;\\
    0.755\,\mathrm{years}\ &\ \dots\dots\dots\dots\ 1435\,\mathrm{mm}^2;\\
    T &= 15.026\,\mathrm{years};
  \end{align*}
  Displacement $6 \rightarrow 7$
  \begin{align*}
    \Delta S_{67} &\approx \frac{203\,\mathrm{mm}\times30\,\mathrm{mm}}{2}
      + 120\,\mathrm{mm}^2 \approx 3165\,\mathrm{mm}^2;\\
    t_6 &= 2001.169\,\mathrm{years};\; t_7 = 2002.831\,\mathrm{years};\\
    \Delta t_{67} &= 1.662\,\mathrm{years};\\
    T\ &\ \dots\dots\dots\dots\ 28560\,\mathrm{mm}^2;\\
    1.662\,\mathrm{years}\ &\ \dots\dots\dots\dots\ 3165\,\mathrm{mm}^2;\\
    T &= 14.997\,\mathrm{years};
  \end{align*}
  Displacement $7 \rightarrow 8$
  \begin{align*}
    \Delta S_{78} &\approx \frac{143\,\mathrm{mm}\times20\,\mathrm{mm}}{2}
      + 5\,\mathrm{mm}^2 \approx 1435\,\mathrm{mm}^2;\\
    t_7 &= 2002.831\,\mathrm{years};\; t_8 = 2003.584\,\mathrm{years};\\
    \Delta t_{78} &= 0.753\,\mathrm{years};\\
    T\ &\ \dots\dots\dots\dots\ 28560\,\mathrm{mm}^2;\\
    0.753\,\mathrm{years}\ &\ \dots\dots\dots\dots\ 1435\,\mathrm{mm}^2;\\
    T &= 14.986\,\mathrm{years};
  \end{align*}
  Displacement $8 \rightarrow 9$
  \begin{align*}
    \Delta S_{89} &\approx \frac{100\,\mathrm{mm}\times22\,\mathrm{mm}}{2}
      + 7\,\mathrm{mm}^2 \approx 1107\,\mathrm{mm}^2;\\
    t_8 &= 2003.584\,\mathrm{years};\; t_9 = 2004.165\,\mathrm{years};\\
    \Delta t_{89} &= 0.581\,\mathrm{years};\\
    T\ &\ \dots\dots\dots\dots\ 28560\,\mathrm{mm}^2;\\
    0.581\,\mathrm{years}\ &\ \dots\dots\dots\dots\ 1107\,\mathrm{mm}^2;\\
    T &= 14.989\,\mathrm{years};
  \end{align*}
  Displacement $9 \rightarrow 10$
  \begin{align*}
    \Delta S_{9;10} &\approx \frac{60\,\mathrm{mm}\times20\,\mathrm{mm}}{2}
      + 200\,\mathrm{mm}^2 \approx 800\,\mathrm{mm}^2;\\
    t_9 &= 2004.165\,\mathrm{years};\; t_{10} = 2004.585\,\mathrm{years};\\
    \Delta t_{9,10} &\approx 0.420\,\mathrm{years};\\
    T\ &\ \dots\dots\dots\dots\ 28560\,\mathrm{mm}^2;\\
    0.420\,\mathrm{years}\ &\ \dots\dots\dots\dots\ 800\,\mathrm{mm}^2;\\
    T &= 14.994\,\mathrm{years};
  \end{align*}
  Displacement $10 \rightarrow 11$
  \begin{align*}
    \Delta S_{10;11} &\approx \frac{20\,\mathrm{mm}\times13\,\mathrm{mm}}{2}
      + 5\,\mathrm{mm}^2 \approx 135\,\mathrm{mm}^2;\\
    t_{10} &= 2004.585\,\mathrm{years};\; t_{11} = 2004.655\,\mathrm{years};\\
    \Delta t_{10;11} &\approx 0.070\,\mathrm{years};\\
    T\ &\ \dots\dots\dots\dots\ 28560\,\mathrm{mm}^2;\\
    0.070\,\mathrm{years}\ &\ \dots\dots\dots\dots\ 135\,\mathrm{mm}^2;\\
    T &= 14.808\,\mathrm{years};
  \end{align*}
  Displacement $11 \rightarrow 12$
  \begin{align*}
    \Delta S_{11;12} &= \frac{22\,\mathrm{mm}\times13\,\mathrm{mm}}{2}
      + 4\,\mathrm{mm}^2 = 147\,\mathrm{mm}^2;\\
    t_{11} &= 2004.655\,\mathrm{years};\; t_{12} = 2004.734\,\mathrm{years};\\
    \Delta t_{11;12} &\approx 0.077\,\mathrm{years};\\
    T\ &\ \dots\dots\dots\dots\ 28560\,\mathrm{mm}^2;\\
    0.077\,\mathrm{years};\ &\ \dots\dots\dots\dots\ 147\,\mathrm{mm}^2;\\
    T &= 14.960\,\mathrm{years};
  \end{align*}
  Displacement $12 \rightarrow 13$
  \begin{align*}
    \Delta S_{12;13} &= \frac{30\,\mathrm{mm}\times13\,\mathrm{mm}}{2}
      + 4\,\mathrm{mm}^2 = 200\,\mathrm{mm}^2;\\
    t_{12} &= 2004.734\,\mathrm{years};\; t_{13} = 2004.839\,\mathrm{years};\\
    \Delta t_{12;13} &\approx 0.105\,\mathrm{years};\\
    T\ &\ \dots\dots\dots\dots\ 28560\,\mathrm{mm}^2;\\
    0.105\,\mathrm{years}\ &\ \dots\dots\dots\dots\ 200\,\mathrm{mm}^2;\\
    T &= 14.994\,\mathrm{years};
  \end{align*}
  Displacement $13 \rightarrow 14$
  \begin{align*}
    \Delta S_{13;14} &= \frac{41\,\mathrm{mm}\times9\,\mathrm{mm}}{2}
      \approx 185\,\mathrm{mm}^2;\\
    t_{13} &= 2004.839\,\mathrm{years};\; t_{14} = 2004.936\,\mathrm{years};\\
    \Delta t_{13;14} &\approx 0.097\,\mathrm{years};\\
    T\ &\ \dots\dots\dots\dots\ 28560\,\mathrm{mm}^2;\\
    0.097\,\mathrm{years}\ &\ \dots\dots\dots\dots\ 185\,\mathrm{mm}^2;\\
    T &= 14.974\,\mathrm{years};
  \end{align*}
  Displacement $14 \rightarrow 15$
  \begin{align*}
    \Delta S_{14;15} &= \frac{90\,\mathrm{mm}\times22\,\mathrm{mm}}{2}
      + 90\,\mathrm{mm}^2 = 1080\,\mathrm{mm}^2;\\
    t_{14} &= 2004.936\,\mathrm{years};\; t_{15} = 2005.503\,\mathrm{years};\\
    \Delta t_{14;15} &\approx 0.567\,\mathrm{years};\\
    T\ &\ \dots\dots\dots\dots\ 28560\,\mathrm{mm}^2;\\
    0.567\,\mathrm{years}\ &\ \dots\dots\dots\dots\ 1080\,\mathrm{mm}^2;\\
    T &= 14.994\,\mathrm{years};
  \end{align*}
  Displacement $15 \rightarrow 16$
  \begin{align*}
    \Delta S_{15;16} &= \frac{1}{2}\cdot120\,\mathrm{mm}\times17\,\mathrm{mm}
      + 5\,\mathrm{mm}^2 = 1025\,\mathrm{mm}^2;\\
    t_{15} &= 2005.503\,\mathrm{years};\; t_{16} = 2006.041\,\mathrm{years};\\
    \Delta t_{15;16} &\approx 0.538\,\mathrm{years};\\
    T\ &\ \dots\dots\dots\dots\ 28560\,\mathrm{mm}^2;\\
    0.538\,\mathrm{years}\ &\ \dots\dots\dots\dots\ 1025\,\mathrm{mm}^2;\\
    T &= 14.990\,\mathrm{years};
  \end{align*}
  Displacement $16 \rightarrow 17$
  \begin{align*}
    \Delta S_{16;17} &= \frac{168\,\mathrm{mm}\times23\,\mathrm{mm}}{2}
      + 10\,\mathrm{mm}^2 = 1942\,\mathrm{mm}^2;\\
    t_{16} &= 2006.041\,\mathrm{years};\; t_{17} = 2007.060\,\mathrm{years};\\
    \Delta t_{16;17} &= 1.019\,\mathrm{years};\\
    T\ &\ \dots\dots\dots\dots\ 28560\,\mathrm{mm}^2;\\
    1.019\,\mathrm{years}\ &\ \dots\dots\dots\dots\ 1942\,\mathrm{mm}^2;\\
    T &= 14.985\,\mathrm{years};
  \end{align*}
  Displacement $17 \rightarrow 1$
  \begin{align*}
    \Delta S_{17;1} &= \frac{240\,\mathrm{mm}\times43\,\mathrm{m}}{2}
      % sic: "43 m" as printed (should be mm)
      + 300\,\mathrm{mm}^2 = 5460\,\mathrm{mm}^2;\\
    t_{17} &= 2007.060\,\mathrm{years};\\
    \Delta t_{17;1} &= 2.867\,\mathrm{years};\\
    T\ &\ \dots\dots\dots\dots\ 28560\,\mathrm{mm}^2;\\
    2.867\,\mathrm{years}\ &\ \dots\dots\dots\dots\ 5460\,\mathrm{mm}^2;\\
    T &= 14.996\,\mathrm{years};\\
    \overline{T} &= 14.923\,\mathrm{years};\\
    T &\approx 15\,\mathrm{years}.
  \end{align*}
\item[c)] By using the III rd Kepler's law:
  \begin{align*}
    T^2 &= \frac{4\pi^2}{K(M_{\mathrm{SA}*} + M_{\mathrm{S}*})}\cdot a^3;\\
    M_{\mathrm{SA}*} + M_{\mathrm{S}*} &= \frac{4\pi^2 a^3}{KT^2};
  \end{align*}
  \[
    a = 4.6125\,\mathrm{light\ days};\;
    T \approx 15\,\mathrm{years};\;
    K = 6.67\cdot10^{-11}\ \mathrm{Nm^2kg^{-2}};
  \]
  \[
    1\,\mathrm{ld} \approx 2.6\cdot10^{13}\,\mathrm{m};\;
    1\,\mathrm{years} \approx 3.1\cdot10^{7}\,\mathrm{s};\;
    4\pi^2 \approx 40;
  \]
  \[
    (2.6)^3 \approx 17.6;\quad (4.6125)^3 \approx 10^2;
  \]
  \begin{align*}
    M_{\mathrm{SA}*} + M_{\mathrm{S}*}
      &= \frac{40\cdot10^{2}\cdot17.6\cdot10^{39}}
              {6.67\cdot10^{-11}\cdot225\cdot10^{15}}\ \mathrm{kg};\\
    M_{\mathrm{SA}*} + M_{\mathrm{S}*} &\approx 4.6\cdot10^{36}\,\mathrm{kg};\\
    M_{\mathrm{Sun}} &\approx 1.8\cdot10^{30}\,\mathrm{kg};
    % sic: as printed (the usual value is 2.0e30 kg)
    \\
    M_{\mathrm{SA}*} + M_{\mathrm{S}*}
      &\approx 2.5\cdot10^{6}\cdot M_{\mathrm{Sun}};\\
    M_{\mathrm{S}*} &\ll M_{\mathrm{SA}*};\\
    M_{\mathrm{SA}*} &\approx 2.5\cdot10^{6}\cdot M_{\mathrm{Sun}},
  \end{align*}
\end{description}

\solhead{11.14}{Hypervelocity stars}
% source id: 2020_DA_3
\begin{description}
\item[(a)] By comparing with spectral charts in appendix 3, we notice that
  \begin{itemize}
  \item H-lines are the most promienent % sic
    lines,
  \item He I feature at 4144\,\AA\ is bearly % sic
    visible and
  \item Ca I feature around 4230\,\AA\ is not seen.
  \end{itemize}
  Thus, it may be reasonable to restrict the range of spectral types to
  B5--A5. \hfill(2.0)

  We find that the ratio of the intensity of a pair of lines and compare it
  with the standard spectra. We see,
  Ca II/H$\epsilon$ $\sim \frac{1}{4}$ \hfill(2.0)

  For B8 standard spectra it is $\sim \frac{1}{5}$, for A0 standard spectra it
  is $\sim \frac{1}{3}$. \hfill(2.0)

  Thus, we deduce that the spectral type lies somewhere between B8 and A0.
  One may write any spectral type in the range, say B9. \hfill(1.0)
\item[(b)] \mbox{}
  \begin{description}
  \item[i.] Appendix 1 gives absolute magnitude for stars of different
    spectral and luminosity classes. The relevant rows are A0 and B5.
    \hfill(1.0)

    In order to estimate the absolute magnitude of a B9 star, interpolate
    between B5 and A0.

    For MS:
    \begin{align*}
      M(\mathrm{B9}) &= M(\mathrm{A0})
        - \frac{(M(\mathrm{A0}) - M(\mathrm{B5}))}{5}\times(10-9)\\
        &= 0.7 - \frac{(0.7 - (-1.1))}{5} = 0.7 - 0.36\\
        &= 0.34 \tag*{(2.0)}
    \end{align*}
    For subgiants:
    \begin{align*}
      M(\mathrm{B9}) &= M(\mathrm{A0})
        - \frac{(M(\mathrm{A0}) - M(\mathrm{B5}))}{5}\times(10-9)\\
        &= 0.1 - \frac{(0.1 - (-1.8))}{5} = 0.1 - 0.38\\
        &= -0.28 \tag*{(2.0)}
    \end{align*}
  \item[ii.]
    \begin{align*}
      m_v &= M_v + 5\log(d) - 5\\
      \therefore d &= 10^{\frac{m_v - M_v + 5}{5}} \tag*{(1.0)}
    \end{align*}
    For MS:
    \[
      d \approx 79.4\,\mathrm{kpc}
    \]
    \hfill(1.0)

    For subgiants:
    \[
      d \approx 105.7\,\mathrm{kpc}
    \]
    \hfill(1.0)
  \item[iii.] \textbf{Yes}. HSV1 % sic
    is off the galactic plane, well into the halo. interstellar absorption
    should not be very high. Thus, ignoring it will not affect the distance
    very much. \hfill(2.0)
  \item[iv.] (A) if it is a MS star,
    \[
      \pi = 1/d \sim 0.01\,\mathrm{mas} < \pi_{\mathrm{GAIA}}
    \]
    Hence, the answer is ``\textbf{No}''. \hfill(1.0)

    (B) If it is a subgiant, then the distance is even larger. Hence, the
    answer is again ``\textbf{No}''. \hfill(1.0)
  \item[v.]
    \ioaafig{2020_DA_3-s1.pdf}
    The left panel is to the scale. The right panel shows details of the
    drawing. \hfill(2.0)
    \begin{align*}
      d(HC)^2 &= d(HP)^2 + d(PC)^2\\
        &= d(HP)^2 + \left[d(SC)^2 + d(PS)^2
           - 2\cdot d(SC)\cdot d(PS)\cos(\angle PSC)\right]\\
      r^2 &= (d\sin b)^2 + r_\odot^2 + (d\cos b)^2
           - 2r_\odot d\cos b\cos(360^\circ - l)\\
      \therefore r &= \sqrt{d^2 + r_\odot^2 - 2r_\odot d\cos b\cos l}\\
        &= \sqrt{105.7^2 + 8^2
           - 2\times8\times105.7\cos(31.332^\circ)\cos(227.335^\circ)}\\
      \therefore r &\approx 110\,\mathrm{kpc} \tag*{(4.0)}
    \end{align*}
  \end{description}
\item[(c)] \mbox{}
  \begin{description}
  \item[i.] From the spectrum, we can estimate the line shift as
    \[
      \Delta\lambda = (11\pm2)\,\text{\AA},\ \lambda_0 = 3934\,\text{\AA}
    \]
    \hfill(3.0)

    Doppler shift gives the heliocentric velocity
    \[
      v_{\mathrm{rHC}} = \frac{c\Delta\lambda}{\lambda_0}
        = (840\pm150)\,\mathrm{km\,s^{-1}}
    \]
    \hfill(2.0)
  \item[ii.]
    \begin{align*}
      v_{\mathrm{rf}} &= v_{\mathrm{rHC}} + 11.1\cos l\cos b
          + 247.24\sin l\,\cos b + 7.25\sin b\\
        &= 840 + 11.1\cos(227.335^\circ)\cos(31.332^\circ)\\
        &\quad + 247.24\sin(227.335^\circ)\cos(31.332^\circ)
          + 7.25\sin(31.332^\circ)\\
      v_{\mathrm{rf}} &= (680\pm130)\,\mathrm{km\,s^{-1}} \tag*{(2.0)}
    \end{align*}
  \item[iii.] The length of an arc does not depend on the coordinate system.
    Thus,
    \begin{align*}
      \Delta\mu &= \sqrt{\mu_\alpha^2 + \mu_\delta^2}\\
        &= \sqrt{0.08^2 + 0.12^2}\\
      \Delta\mu &= (0.144\pm0.340)\,\mathrm{mas/yr} \tag*{(2.0)}\\
      v_\theta = d\Delta\mu
        &= \frac{105.7\times10^{3}\times3.086\times10^{13}
           \times0.144\times10^{-3}}{206\,265\times3.16\times10^{7}}\\
        &= (72\pm170)\,\mathrm{km\,s^{-1}} \tag*{(2.0)}
    \end{align*}
  \item[iv.] The direction to the Galactic center is \emph{also}
    perpendicular to the plane of the sky. Thus,
    \begin{align*}
      v_{\mathrm{T}} &= \sqrt{v_{\mathrm{rf}}^2 + v_\theta^2}
        = \sqrt{680^2 + 72^2}\\
      v_{\mathrm{T}} &= (684\pm210)\,\mathrm{km\,s^{-1}} \tag*{(2.0)}
    \end{align*}
    i.e.\ it is almost radial.
    \[
      \theta = \tan^{-1}\left(\frac{v_\theta}{v_{\mathrm{rf}}}\right)
        \approx \tan^{-1}(0.105) \approx 6^\circ
    \]
    \hfill(2.0)
  \item[v.] As the velocity is almost radial, extrapolating backwards, the
    star would be in galactic disk somewhere close to the galactic centre.
    Thus, the answer is \textbf{A}. \hfill(2.0)
  \end{description}
\item[(d)] \mbox{}
  \begin{description}
  \item[i.]
    \[
      v_{esc} = \sqrt{\frac{2GM}{r}}
    \]
    \hfill(1.0)
  \item[ii.]
    \begin{align*}
      M_r &= 4\pi\rho_0 r_c^2
        \left[r - r_c\tan^{-1}\left(\frac{r}{r_c}\right)\right]\\
        &= 4\pi\times1.396\times10^{4}\times8^2
        \left[110 - 8\times\tan^{-1}\left(\frac{110}{8}\right)\right]\\
      M_r(110\,\mathrm{kpc}) &\approx 1.2\times10^{12}M_\odot \tag*{(2.0)}
    \end{align*}
  \item[iii.]
    \begin{align*}
      v_{esc} &= \sqrt{\frac{2GM}{r}}
        = \sqrt{\frac{2\times6.674\times10^{-11}\times1.2\times10^{12}
          \times1.998\times10^{30}}{110\times10^{3}\times3.086\times10^{16}}}\\
        &= 311\,\mathrm{km\,s^{-1}} \tag*{(2.0)}
    \end{align*}
  \item[iv.] $v_{\mathrm{rf}} > v_{esc} \implies$ \textbf{Yes}. \hfill(1.0)
  \end{description}
\item[(e)] If HSV1 % sic
  originated close to the galactic center and assuming is was % sic
  ejected at the measured speed, it has been traveling for
  \[
    t_{\mathrm{trav}} = \frac{r}{v_{\mathrm{rf}}}
      = \frac{110\times10^{3}\times3.086\times10^{16}}{680\times10^{3}}
      = 1.6\times10^{8}\,\mathrm{yr}
  \]
  \hfill(2.0)
\item[(f)] Main Sequence life time for the star will approximately be
  \[
    t \sim \frac{M}{L} \sim L^{-5/7}
  \]
  Comparing with absolute magnitude of the Sun,
  \[
    t \sim 10^{8}\,\mathrm{yr}
  \]
  This is very similar to the time taken by the star to travel from the disk
  to present position. Thus, the anwer % sic
  is \textbf{A}. \hfill(3.0)
\item[(g)] Stars of early type are very young, so they cannot belong to the
  native population of the halo. They must be coming from elsewhere. As a
  result, they may be moving very fast and some of them may be on their way
  to abandoning the Milky Way.
  $\implies$ C. \hfill(2.0)
\end{description}

\solhead{11.15}{Shapley Hypothesis}
% source id: 2024_DA_D2
\begin{description}
\item[(a)] The first step is to convert the extinction corrected distance
  moduli ($DM$) of all globular clusters to distance ($d$), in parsecs:
  \[
    DM = (m - M) - Av = 5\cdot\log(d) - 5
    \quad\Rightarrow\quad
    d = 10^{(DM+5)/5}\,\mathrm{pc}
  \]
  So, it follows to calculate the cartesian coordinates ($x$, $y$, $z$) of the
  globular clusters with respect to the Sun, using the galactic coordinates
  (longitude l and latitude b). The $x$-axis points to the Galactic Centre,
  the $y$-axis points to direction of galactic rotation and the $z$-axis is
  perpendicular to the galactic disk, and points in the direction antiparallel
  to the angular momentum. The conversion of coordinates should be done as
  follows:
  \[
    x = d\cdot\cos(l)\cdot\cos(b)\;,\quad
    y = d\cdot\sin(l)\cdot\cos(b)\;,\quad
    z = d\cdot\sin(b)
  \]
  The table below shows the calculated values of $d$, $x$, $y$, and $z$ for
  all globular clusters.

  \begin{center}
  \begin{tabular}{|l|r|r|r|r|}
  \hline
  \textbf{Name} & \textbf{d (pc)} & \textbf{x (pc)} & \textbf{y (pc)} & \textbf{z (pc)}\\
  \hline
  NGC 6522    &  7244.36 &  7226.20 &   129.29 &  -496.01\\
  NGC 6401    &  7585.78 &  7553.77 &   455.39 &   526.52\\
  NGC 6342    &  7943.28 &  7800.55 &   668.47 &  1341.77\\
  NGC 6553    &  5248.07 &  5218.73 &   479.81 &  -277.32\\
  NGC 6440    &  8317.64 &  8223.94 &  1116.16 &   551.39\\
  Ter 12      &  5248.07 &  5188.85 &   762.34 &  -192.40\\
  VVV-CL 160  &  6918.31 &  6809.92 &  1219.29 &    36.47\\
  2MASS-GC01  &  3311.31 &  3256.16 &   601.79 &     5.78\\
  NGC 6517    &  9120.11 &  8551.60 &  2982.17 &  1073.85\\
  NGC 6402    &  9120.11 &  8213.73 &  3206.36 &  2330.31\\
  NGC 6712    &  7244.36 &  6528.00 &  3093.30 &  -545.44\\
  NGC 6426    & 20892.96 & 17697.53 &  9444.44 &  5840.86\\
  NGC 5466    & 15848.93 &  3319.16 &  3004.35 & 15203.48\\
  NGC 7089    & 11481.54 &  5558.08 &  7476.05 & -6711.34\\
  NGC 288     &  9120.11 &   -86.55 &    47.41 & -9119.57\\
  NGC 2298    & 10000.00 & -3966.46 & -8755.80 & -2757.38\\
  NGC 4590    & 10471.29 &  4185.03 & -7359.22 &  6162.41\\
  NGC 4372    &  5754.40 &  2919.15 & -4859.63 &  -987.77\\
  NGC 362     &  8709.64 &  3150.05 & -5133.77 & -6291.21\\
  BH 140      &  4786.30 &  2611.38 & -3995.02 &  -359.45\\
  NGC 5927    &  8317.64 &  6919.31 & -4561.75 &   704.68\\
  Patchick 126&  7943.28 &  7465.49 & -2661.12 &  -530.03\\
  NGC 5897    & 12589.25 & 10392.20 & -3187.93 &  6350.49\\
  NGC 6380    &  9549.93 &  9393.28 & -1625.54 &  -570.03\\
  Djor 1      & 10000.00 &  9973.79 &  -579.45 &  -433.40\\
  \hline
  \end{tabular}
  \end{center}
\item[(b)] In order to estimate the distance from the Sun to the centre of
  the distribution and its uncertainty, we first need the mean values of each
  coordinate $\bar{x}$, $\bar{y}$ and $\bar{z}$, as well as the standard
  deviations in each axis, $\sigma_x$, $\sigma_y$ and $\sigma_z$. The
  calculations of the mean and standard deviation for the $x$, $y$, and $z$
  coordinates are as follows:

  \textbf{Mean in the three axes:}
  \[
    \bar{x} = \frac{1}{N}\sum_{i=1}^{N}x_i,\quad
    \bar{y} = \frac{1}{N}\sum_{i=1}^{N}y_i,\quad
    \bar{z} = \frac{1}{N}\sum_{i=1}^{N}z_i
  \]
  \[
    \bar{x} = 6164.1\ \mathrm{pc},\quad
    \bar{y} = -321.3\ \mathrm{pc},\quad
    \bar{z} = 434.3\ \mathrm{pc}
  \]

  \textbf{Standard deviations in the three axes:}
  \[
    \sigma_x = \sqrt{\frac{\sum_{i=1}^{N}(x_i-\bar{x})^2}{N-1}},\quad
    \sigma_y = \sqrt{\frac{\sum_{i=1}^{N}(y_i-\bar{y})^2}{N-1}},\quad
    \sigma_z = \sqrt{\frac{\sum_{i=1}^{N}(z_i-\bar{z})^2}{N-1}}
  \]
  \[
    \sigma_x = 4057.7\,\mathrm{pc}\;,\quad
    \sigma_y = 4196.4\,\mathrm{pc}\;,\quad
    \sigma_z = 4682.6\,\mathrm{pc}
  \]
  Consequently, we obtain the uncertainties associated to the mean in each
  axis:

  \textbf{Uncertainties of means in the three axes:}
  \[
    \delta\bar{x} = \frac{\sigma_x}{\sqrt{N}},\quad
    \delta\bar{y} = \frac{\sigma_y}{\sqrt{N}},\quad
    \delta\bar{z} = \frac{\sigma_z}{\sqrt{N}}
  \]
  \[
    \delta\bar{x} = 811.54\ \mathrm{pc},\quad
    \delta\bar{y} = 839.28\ \mathrm{pc},\quad
    \delta\bar{z} = 936.52\ \mathrm{pc}
  \]
  Finally, the distance $D$, in parsecs, from the Sun to the centre of the
  distribution of globular clusters is estimated as:
  \begin{align*}
    D &= \sqrt{\bar{x}^2 + \bar{y}^2 + \bar{z}^2}\\
    D &= \sqrt{(6164.1)^2 + (-321.3)^2 + (434.3)^2} = 6187.73\ \mathrm{pc}
  \end{align*}
  And, to estimate its uncertainty $\delta D$, we should perform an error
  propagation calculation:
  \begin{align*}
    (\delta D)^2 &= \left(\frac{\partial D}{\partial\bar{x}}\right)^2
        \cdot(\delta\bar{x})^2
      + \left(\frac{\partial D}{\partial\bar{y}}\right)^2
        \cdot(\delta\bar{y})^2
      + \left(\frac{\partial D}{\partial\bar{z}}\right)^2
        \cdot(\delta\bar{z})^2 \Rightarrow\\
    (\delta D)^2 &= \left[
        \left(\frac{\bar{x}}{D}\cdot\delta\bar{x}\right)^2
      + \left(\frac{\bar{y}}{D}\cdot\delta\bar{y}\right)^2
      + \left(\frac{\bar{z}}{D}\cdot\delta\bar{z}\right)^2
      \right] \Rightarrow\\
    \delta D &= \frac{1}{6187.73}\cdot
      \left[(6164.1\cdot811.54)^2 + (-321.3\cdot839.28)^2
      + (434.3\cdot936.52)^2\right]^{1/2} \Rightarrow\\
    \delta D &= 812.28\,\mathrm{pc}
  \end{align*}
\item[(c)] The table below shows the calculated values of the quartiles for
  each distribution.

  \begin{center}
  \begin{tabular}{|c|r|r|r|}
  \hline
  \textbf{Axis} & \textbf{Q1 (pc)} & \textbf{Q2 (pc)} & \textbf{Q3 (pc)}\\
  \hline
  $x$ &  3287.66 & 6809.92 & 8218.84\\
  $y$ & -3591.48 &  455.39 & 2100.73\\
  $z$ &  -557.74 & -192.40 & 1207.81\\
  \hline
  \end{tabular}
  \end{center}

  The expected histograms for each axis are as follows:
  \ioaafig{2024_DA_D2-s1-x.pdf}
  \ioaafig{2024_DA_D2-s1-y.pdf}
  \ioaafig{2024_DA_D2-s1-z.pdf}
\item[(d)] \textbf{F} (False). The analysed sample does not follow Shapley's
  hypothesis, because all distributions are asymmetrical. See table below for
  the values of calculated symmetry factors for each distribution.

  \begin{center}
  \begin{tabular}{|c|c|l|}
  \hline
  \textbf{Axis} & $\Phi$ & \textbf{Symmetry Type}\\
  \hline
  $x$ & 0,429 & asymmetrical\\   $y$ & 0,422 & asymmetrical\\
  $z$ & 0,588 & asymmetrical\\
  \hline
  \end{tabular}
  \end{center}
\end{description}

\solhead{11.16}{Identify open/globular clusters and nebulae}
% source id: 2012_OBS-TEL_4
Targets pointed by the evaluator will be M22, M8, M7 and M57. Student will
have 30 seconds to evaluate each target.

\begin{description}
\item[Object 1 (M22)] (GC)
\item[Object 2 (M57)] (PN)
\item[Object 3 (M7)] (OC)
\item[Object 4 (M8)] (EN)
\end{description}
